\fbox{例題 2.1}

\textbf{手順1} 測定単位$m = 5$、データ数$n = 50$

\textbf{手順2} データの最大値$x_{\text{max}} = 36.5$、データの最小値$x_{\text{min}} = 11.0$、範囲$R = 36.5 - 11.0 = 25.5$

\textbf{手順3} 仮の区間の数$h \doteqdot \sqrt{n} = \sqrt{50} \rightarrow 7$

\textbf{手順4} 区間の幅$c = \dfrac{R}{h} = \dfrac{25.5}{7} = 3.64 \rightarrow 3.5$

\textbf{手順5} 一番下の境界値を$x_{\text{min}} - \dfrac{m}{2} = 11.0 - \dfrac{0.5}{2} = 10.75$と定め$c = 3.5$を順次加えていく

\textbf{手順6、7} 各区間の中心値を求め、度数をチェックして数え上げる

\vskip\baselineskip
\begin{tabular}{c|c|c} \hline
  区間 & 中心値 & 度数$f$ \\ \hline
  10.75〜14.25 & 12.50 & 2 \\
  14.25〜17.75 & 16.00 & 4 \\
  17.75〜21.25 & 19.50 & 7 \\
  21.25〜24.75 & 23.00 & 14 \\
  24.75〜28.25 & 26.50 & 12 \\
  28.25〜31.75 & 30.00 & 6 \\
  31.75〜35.25 & 33.50 & 3 \\
  35.25〜38.75 & 37.00 & 2 \\ \hline
\end{tabular}
